\section{Discussion and limitations}
\label{sec:discuss}

\paragraph{What the typed guarantee buys.} The specialist ensemble is the natural engineering baseline, and on raw detection it is competitive or better: Page--Hinkley and the conformal martingales detect concept changes sooner than any TSSD witness. Its \emph{types}, however, are not calibrated. A covariate shift that changes prevalence fires the label-shift CUSUM, and a mixed covariate--label shift fires the loss monitor and is reported as concept. Each of these errors maps to a costly action: an unnecessary prior correction, or an unnecessary retraining. TSSD's concept claim cannot be fooled in this way at more than its nominal rate, whatever the inputs do. Across all our experiments, on 16{,}830 streams with or without covariate or label shift, the claim never fired on a stream whose target posterior was prior-correctable.

\paragraph{The price.} The certified claim pays twice. First, the alternative must beat a null fitted with hindsight that also allows a change point, so concept evidence accumulates more slowly than in an uncorrected loss monitor. Second, the evidence is in labels: with $K=10$ classes and about seven post-change labels (Digits at 5\%), or with natural shifts in which a prior correction already removes most of the excess loss (Covertype), a certified claim is out of reach at a 5\% label budget, and TSSD correctly does not make it. With four times the label budget the Digits concept claims become reliable. This information bottleneck is intrinsic to the claim, not an artefact of the implementation: no procedure can certify ``not prior-correctable'' faster than the labels reveal it.

\paragraph{Covariate versus label.} The distinction between these two types remains heuristic. A covariate shift whose labels follow the source posterior generally changes prevalence, so ``no label shift'' is not identifiable from inputs alone without a model of $P(X\mid Y)$. TSSD's orthogonalisation and predictable gate make the distinction work in our controlled and natural settings, but the error rates in Section~\ref{sec:typing} are empirical. The action consequence of this residual error is mild, since both types leave the posterior prior-correctable.

\paragraph{Assumptions.} The guarantees require random label reveal (A1). Under active or error-driven labelling the concept claim is not protected. The global-null statements for the $C$ witness also require the source posterior to lie in the tempered prior-correction family (A3). The tempering grid gives some slack, and our null and same-domain streams showed no false concept claims, but a badly miscalibrated source model would need recalibration before deployment. Time-series dependence, as in Electricity, violates the exchangeability that underlies the paired witnesses. Those segments are therefore descriptive rather than tests of validity, and the null-period check there is only indicative. Finally, $H^C$ allows one change of the prior-correction parameter per SR start. A drifting prior is covered only approximately, through the SR restarts.

\paragraph{Cost.} A 400-step stream takes 1--2\,s on one CPU core in pure Python. The cost of the $C$ witness grows quadratically in the number of revealed labels within an SR start, which we cap at 40 starts. For high label rates or very long streams, windowing the suffix computation, or replacing it with a sufficient-statistic update, is an obvious engineering improvement.

\section{Conclusion}
\label{sec:conclusion}

TSSD turns three sources of evidence (unlabelled inputs, prevalence and sparse labels) into a single monitor with an anytime-valid alarm, a restarting change-detection mode with average-run-length control, and a concept claim certified against covariate and label shift. Across controlled, semi-real and naturally shifted data, and against ten detectors from the recent literature, the experiments support a precise claim. TSSD is not the fastest detector of each shift type, but it is the only one among them whose typed output is reliable enough to drive the keep/recalibrate/retrain decision. Natural next steps are an analogous guarantee for the covariate/label distinction under explicit class-conditional models, and label-efficient alternatives for large label spaces.
